\section*{习题 6.5\quad 广义积分}

1. 判断题

(1) \(\int_{-1}^{1}\frac{1}{x}\,dx=0\) 是否正确？

解：不正确。因为被积函数 \(\frac{1}{x}\) 在 \(x=0\) 处无界，且积分区间包含奇点 \(x=0\)。该积分是瑕积分，需要分别计算：
\[ \int_{-1}^{1}\frac{1}{x}\,dx=\lim_{\epsilon\to0^{+}}\int_{-1}^{-\epsilon}\frac{1}{x}\,dx+\lim_{\delta\to0^{+}}\int_{\delta}^{1}\frac{1}{x}\,dx. \]
由于 \(\int\frac{1}{x}\,dx=\ln|x|+C\)，所以：
\[ \lim_{\epsilon\to0^{+}}\big[\ln|x|\big]_{-1}^{-\epsilon}=\lim_{\epsilon\to0^{+}}(\ln\epsilon-\ln1)=-\infty, \]
\[ \lim_{\delta\to0^{+}}\big[\ln|x|\big]_{\delta}^{1}=\lim_{\delta\to0^{+}}(\ln1-\ln\delta)=-\infty. \]
两个部分都发散，因此原积分发散。故该等式不正确。

(2) \(\int_{-\infty}^{+\infty}\frac{1}{x^{2}}\,dx=0\) 是否正确？

解：不正确。因为被积函数 \(\frac{1}{x^{2}}\) 在 \(x=0\) 处无界，且积分区间包含奇点 \(x=0\)。该积分需要分别计算：
\[ \int_{-\infty}^{+\infty}\frac{1}{x^{2}}\,dx=\int_{-\infty}^{0}\frac{1}{x^{2}}\,dx+\int_{0}^{+\infty}\frac{1}{x^{2}}\,dx. \]
由于 \(\int\frac{1}{x^{2}}dx=-\frac{1}{x}+C\)，所以：
\[ \int_{-\infty}^{0}\frac{1}{x^{2}}dx=\lim_{a\to-\infty}\int_{a}^{0}\frac{1}{x^{2}}dx=\lim_{a\to-\infty}\left[-\frac{1}{x}\right]_{a}^{0}=\lim_{a\to-\infty}\left(-\frac{1}{0^{+}}+\frac{1}{a}\right)=-\infty. \]
类似地，\(\int_{0}^{+\infty}\frac{1}{x^{2}}dx\) 也发散。因此原积分发散，故该等式不正确。

(3) \(\int_{-\infty}^{+\infty}\sin x\,dx=\lim_{a\to+\infty}\int_{-a}^{a}\sin x\,dx=\lim_{a\to+\infty}0=0\) 是否正确？

解：不正确。虽然 \(\int_{-a}^{a}\sin x\,dx=0\) 对于任意有限 \(a\) 成立，但广义积分 \(\int_{-\infty}^{+\infty}\sin x\,dx\) 的定义是：
\[ \int_{-\infty}^{+\infty}\sin x\,dx=\lim_{a\to-\infty,\,b\to+\infty}\int_{a}^{b}\sin x\,dx, \]
而不仅仅是对称极限 \(\lim_{a\to+\infty}\int_{-a}^{a}\sin x\,dx\)。如果取非对称极限，例如：
\[ \lim_{a\to-\infty,\,b\to+\infty}\int_{a}^{b}\sin x\,dx=\lim_{a\to-\infty,\,b\to+\infty}(-\cos b+\cos a), \]
该极限不存在。因此，广义积分 \(\int_{-\infty}^{+\infty}\sin x\,dx\) 发散。故该等式不正确。

(4) 当 \(p<1\) 时，\(\int_{0}^{2}\frac{1}{x^{p}(1-x)}\,dx\) 收敛是否正确？

解：不正确。该积分在 \(x=0\) 和 \(x=1\) 处都可能为奇点。在 \(x=0\) 处，被积函数行为类似于 \(\frac{1}{x^{p}}\)，当 \(p<1\) 时收敛；但在 \(x=1\) 处，被积函数行为类似于 \(\frac{1}{1-x}\)，而 \(\int\frac{1}{1-x}dx\) 在 \(x=1\) 处发散。因此，即使 \(p<1\)，积分在 \(x=1\) 处发散，故整个积分发散。该陈述不正确。

(5) 对任意的 \(p>0\)，\(\int_{0}^{+\infty}\frac{1}{x^{p}}dx\) 均发散是否正确？

解：该积分在 \(x=0\) 和 \(x\to+\infty\) 处都需要考虑。
\begin{itemize}
\item 当 \(p>1\) 时，\(\int_{1}^{+\infty}\frac{1}{x^{p}}dx\) 收敛，但 \(\int_{0}^{1}\frac{1}{x^{p}}dx\) 发散（因为 \(p>1\)）。
\item 当 \(0<p<1\) 时，\(\int_{0}^{1}\frac{1}{x^{p}}dx\) 收敛，但 \(\int_{1}^{+\infty}\frac{1}{x^{p}}dx\) 发散。
\item 当 \(p=1\) 时，两个部分都发散。
\end{itemize}
因此，对于任意 \(p>0\)，该积分总是发散。故该陈述正确。

2. 填空题

(1) \(\int_{1}^{+\infty}\frac{\ln x}{x^{2}}\,dx=\)

解答过程：

这是一个广义积分，使用分部积分法求解。令 \(u=\ln x,\ dv=\frac{1}{x^{2}}dx\)，则
\[ du=\frac{1}{x}dx,\qquad v=-\frac{1}{x}. \]
分部积分公式 \(\int u\,dv=uv-\int v\,du\)，代入得：
\[ \int\frac{\ln x}{x^{2}}dx=-\frac{\ln x}{x}-\int\left(-\frac{1}{x}\right)\cdot\frac{1}{x}dx=-\frac{\ln x}{x}+\int\frac{1}{x^{2}}dx=-\frac{\ln x}{x}-\frac{1}{x}+C. \]
计算广义积分：
\[ \int_{1}^{+\infty}\frac{\ln x}{x^{2}}dx=\lim_{b\to+\infty}\left[-\frac{\ln x}{x}-\frac{1}{x}\right]_{1}^{b}=\lim_{b\to+\infty}\left(-\frac{\ln b}{b}-\frac{1}{b}\right)-\left(0-1\right). \]
使用洛必达法则求极限：
\[ \lim_{b\to+\infty}\frac{\ln b}{b}=\lim_{b\to+\infty}\frac{1/b}{1}=0,\qquad \lim_{b\to+\infty}\frac{1}{b}=0. \]
所以 \(\int_{1}^{+\infty}\frac{\ln x}{x^{2}}dx=0-0+1=1\)。

答案：\(1\)

(2) \(\int_{0}^{1}\frac{1}{\sqrt{1-x^{2}}}\,dx=\)

解答过程：

这是一个基本积分公式，可以直接求解：
\[ \int\frac{1}{\sqrt{1-x^{2}}}dx=\arcsin x+C. \]
计算定积分：
\[ \int_{0}^{1}\frac{1}{\sqrt{1-x^{2}}}dx=\big[\arcsin x\big]_{0}^{1}=\arcsin1-\arcsin0=\frac{\pi}{2}-0=\frac{\pi}{2}. \]
答案：\(\frac{\pi}{2}\)

(3) \(\int_{0}^{1}\frac{1}{\sqrt{x}}e^{\sqrt{x}}dx=\)

解答过程：

使用换元法求解。令 \(t=\sqrt{x}\)，则
\[ x=t^{2},\qquad dx=2t\,dt. \]
当 \(x=0\) 时 \(t=0\)；当 \(x=1\) 时 \(t=1\)。代入原积分：
\[ \int_{0}^{1}\frac{1}{\sqrt{x}}e^{\sqrt{x}}dx=\int_{0}^{1}\frac{1}{t}e^{t}\cdot2t\,dt=2\int_{0}^{1}e^{t}dt. \]
计算积分：
\[ 2\int_{0}^{1}e^{t}dt=2\big[e^{t}\big]_{0}^{1}=2(e-1). \]
答案：\(2(e-1)\)

(4) 如果 \(a>0\)，则 \(\int_{-\infty}^{+\infty}e^{-a|x|}dx=\)

解答过程：

由于被积函数是偶函数，可以简化计算：
\[ \int_{-\infty}^{+\infty}e^{-a|x|}dx=2\int_{0}^{+\infty}e^{-ax}dx. \]
计算积分：
\[ \int_{0}^{+\infty}e^{-ax}dx=\lim_{b\to+\infty}\int_{0}^{b}e^{-ax}dx=\lim_{b\to+\infty}\left[-\frac{1}{a}e^{-ax}\right]_{0}^{b}=\lim_{b\to+\infty}\left(-\frac{1}{a}e^{-ab}+\frac{1}{a}\right)=0+\frac{1}{a}=\frac{1}{a}. \]
所以 \(\int_{-\infty}^{+\infty}e^{-a|x|}dx=2\cdot\frac{1}{a}=\frac{2}{a}\)。

答案：\(\frac{2}{a}\)

3. 计算下列广义积分

(1) \(\int_{-\infty}^{+\infty}\frac{dx}{x^{2}+x+1}\)

解：分母配方：\(x^{2}+x+1=\left(x+\frac{1}{2}\right)^{2}+\frac{3}{4}\)。令 \(t=x+\frac{1}{2}\)，则
\[ \int_{-\infty}^{+\infty}\frac{dx}{x^{2}+x+1}=\int_{-\infty}^{+\infty}\frac{dt}{t^{2}+\left(\frac{\sqrt3}{2}\right)^{2}}=\frac{2}{\sqrt3}\left[\arctan\frac{2t}{\sqrt3}\right]_{-\infty}^{+\infty}=\frac{2}{\sqrt3}\left(\frac{\pi}{2}-\left(-\frac{\pi}{2}\right)\right)=\frac{2\pi}{\sqrt3}. \]

(2) \(\int_{2}^{+\infty}\frac{dx}{x^{2}+x-2}\)

解：分母因式分解：\(x^{2}+x-2=(x+2)(x-1)\)。部分分式分解：
\[ \frac{1}{x^{2}+x-2}=\frac{A}{x+2}+\frac{B}{x-1}. \]
解得 \(A=-\frac{1}{3},\ B=\frac{1}{3}\)。所以
\[ \int\frac{dx}{x^{2}+x-2}=-\frac{1}{3}\ln|x+2|+\frac{1}{3}\ln|x-1|+C=\frac{1}{3}\ln\left|\frac{x-1}{x+2}\right|+C. \]
计算广义积分：
\[ \int_{2}^{+\infty}\frac{dx}{x^{2}+x-2}=\lim_{b\to+\infty}\left[\frac{1}{3}\ln\left|\frac{x-1}{x+2}\right|\right]_{2}^{b}=\lim_{b\to+\infty}\frac{1}{3}\left(\ln\left|\frac{b-1}{b+2}\right|-\ln\frac{1}{4}\right)=\frac{1}{3}(\ln1+\ln4)=\frac{\ln4}{3}. \]

(3) \(\int_{0}^{1}x^{2}\ln x\,dx\)

解：令 \(u=\ln x,\ dv=x^{2}dx\)，则 \(du=\frac{1}{x}dx,\ v=\frac{x^{3}}{3}\)。分部积分：
\[ \int x^{2}\ln x\,dx=\frac{x^{3}}{3}\ln x-\int\frac{x^{3}}{3}\cdot\frac{1}{x}dx=\frac{x^{3}}{3}\ln x-\frac{1}{3}\int x^{2}dx=\frac{x^{3}}{3}\ln x-\frac{x^{3}}{9}+C. \]
计算广义积分（\(x=0\) 为奇点）：
\[ \int_{0}^{1}x^{2}\ln x\,dx=\lim_{\epsilon\to0^{+}}\int_{\epsilon}^{1}x^{2}\ln x\,dx=\lim_{\epsilon\to0^{+}}\left[\frac{x^{3}}{3}\ln x-\frac{x^{3}}{9}\right]_{\epsilon}^{1}=\left(0-\frac{1}{9}\right)-\lim_{\epsilon\to0^{+}}\left(\frac{\epsilon^{3}}{3}\ln\epsilon-\frac{\epsilon^{3}}{9}\right)=-\frac{1}{9}-0=-\frac{1}{9}. \]

(4) \(\int_{0}^{+\infty}e^{-\sqrt{x}}dx\)

解：令 \(t=\sqrt{x}\)，则 \(x=t^{2},\ dx=2t\,dt\)。当 \(x=0\) 时 \(t=0\)，当 \(x\to+\infty\) 时 \(t\to+\infty\)。积分变为：
\[ \int_{0}^{+\infty}e^{-\sqrt{x}}dx=\int_{0}^{+\infty}e^{-t}\cdot2t\,dt=2\int_{0}^{+\infty}te^{-t}dt. \]
分部积分：令 \(u=t,\ dv=e^{-t}dt\)，则 \(du=dt,\ v=-e^{-t}\)。所以
\[ \int te^{-t}dt=-te^{-t}+\int e^{-t}dt=-te^{-t}-e^{-t}+C. \]
因此
\[ 2\int_{0}^{+\infty}te^{-t}dt=2\big[-te^{-t}-e^{-t}\big]_{0}^{+\infty}=2\left(0-(0-1)\right)=2. \]

(5) \(\int_{1}^{2}\frac{x}{\sqrt{x^{2}-1}}dx\)

解：令 \(t=\sqrt{x^{2}-1}\)，则 \(t^{2}=x^{2}-1,\ 2t\,dt=2x\,dx\)，即 \(x\,dx=t\,dt\)。当 \(x=1\) 时 \(t=0\)，当 \(x=2\) 时 \(t=\sqrt3\)。积分变为：
\[ \int_{1}^{2}\frac{x}{\sqrt{x^{2}-1}}dx=\int_{0}^{\sqrt3}\frac{t}{t}dt=\int_{0}^{\sqrt3}dt=\big[t\big]_{0}^{\sqrt3}=\sqrt3. \]

(6) \(\int_{0}^{1}\sqrt{\frac{x}{1-x}}\,dx\)

解：令 \(x=\sin^{2}\theta\)，则 \(dx=2\sin\theta\cos\theta\,d\theta\)。当 \(x=0\) 时 \(\theta=0\)，当 \(x=1\) 时 \(\theta=\frac{\pi}{2}\)。积分变为：
\[ \int_{0}^{1}\sqrt{\frac{x}{1-x}}dx=\int_{0}^{\frac{\pi}{2}}\sqrt{\frac{\sin^{2}\theta}{\cos^{2}\theta}}\cdot2\sin\theta\cos\theta\,d\theta=\int_{0}^{\frac{\pi}{2}}2\sin^{2}\theta\,d\theta. \]
利用公式 \(\sin^{2}\theta=\frac{1-\cos2\theta}{2}\)：
\[ \int_{0}^{\frac{\pi}{2}}2\sin^{2}\theta\,d\theta=\int_{0}^{\frac{\pi}{2}}(1-\cos2\theta)d\theta=\left[\theta-\frac{\sin2\theta}{2}\right]_{0}^{\frac{\pi}{2}}=\left(\frac{\pi}{2}-0\right)-(0-0)=\frac{\pi}{2}. \]

(7) \(\int_{a}^{b}\frac{dx}{\sqrt{(x-a)(b-x)}}\)（\(b>a\)）

解：令 \(x=a+(b-a)\sin^{2}\theta\)，则 \(dx=2(b-a)\sin\theta\cos\theta\,d\theta\)。当 \(x=a\) 时 \(\theta=0\)，当 \(x=b\) 时 \(\theta=\frac{\pi}{2}\)。积分变为：
\[ \int_{a}^{b}\frac{dx}{\sqrt{(x-a)(b-x)}}=\int_{0}^{\frac{\pi}{2}}\frac{2(b-a)\sin\theta\cos\theta\,d\theta}{\sqrt{(b-a)\sin^{2}\theta\cdot(b-a)\cos^{2}\theta}}=\int_{0}^{\frac{\pi}{2}}\frac{2(b-a)\sin\theta\cos\theta}{(b-a)\sin\theta\cos\theta}d\theta=\int_{0}^{\frac{\pi}{2}}2\,d\theta=\pi. \]

(8) \(\int_{0}^{1}(\ln x)^{2}dx\)

解：令 \(t=-\ln x\)，则 \(x=e^{-t},\ dx=-e^{-t}dt\)。当 \(x=0\) 时 \(t\to+\infty\)，当 \(x=1\) 时 \(t=0\)。积分变为：
\[ \int_{0}^{1}(\ln x)^{2}dx=\int_{+\infty}^{0}t^{2}\cdot(-e^{-t})dt=\int_{0}^{+\infty}t^{2}e^{-t}dt. \]
这是 Gamma 函数：
\[ \int_{0}^{+\infty}t^{2}e^{-t}dt=\Gamma(3)=2!=2. \]

4. 当 \(k\) 为何值时，广义积分 \(\int_{2}^{+\infty}\frac{dx}{x(\ln x)^{k}}\) 收敛？

解：考虑积分 \(\int_{2}^{+\infty}\frac{dx}{x(\ln x)^{k}}\)。令 \(u=\ln x\)，则 \(du=\frac{1}{x}dx\)。当 \(x=2\) 时 \(u=\ln2\)，当 \(x\to+\infty\) 时 \(u\to+\infty\)。积分变为：
\[ \int_{\ln2}^{+\infty}\frac{du}{u^{k}}. \]
该积分在 \(k>1\) 时收敛，在 \(k\le1\) 时发散。因此，当 \(k>1\) 时原积分收敛。

5. 判别下列积分的敛散性

(1) \(\int_{0}^{+\infty}\frac{x^{2}}{x^{4}-x^{2}+1}dx\)

解：当 \(x\to+\infty\) 时，被积函数 \(\sim\frac{x^{2}}{x^{4}}=\frac{1}{x^{2}}\)，而 \(\int\frac{1}{x^{2}}dx\) 收敛。因此该积分收敛。

(2) \(\int_{0}^{+\infty}\frac{1}{x\sqrt[3]{4x^{2}+1}}dx\)（注：原题可能有误，假设为 \(\int_{0}^{+\infty}\frac{1}{x^{4/3}+1}dx\)）

解：当 \(x\to+\infty\) 时，被积函数 \(\sim\frac{1}{x^{4/3}}\)，而 \(4/3>1\)，收敛。当 \(x\to0\) 时，被积函数 \(\sim1\)，有限。因此积分收敛。

(3) \(\int_{0}^{2}\frac{dx}{\ln x}\)

解：奇点在 \(x=1\) 处。当 \(x\to1\) 时，\(\ln x\sim x-1\)，因此被积函数 \(\sim\frac{1}{x-1}\)，发散。故积分发散。

(4) \(\int_{0}^{+\infty}\frac{x^{m}}{x^{n}+1}dx\)（\(n\ge0\)）

解：需考虑 \(x\to0\) 和 \(x\to+\infty\) 的行为。
\begin{itemize}
\item 当 \(x\to0\) 时，被积函数 \(\sim x^{m}\)，收敛当且仅当 \(m>-1\)。
\item 当 \(x\to+\infty\) 时，被积函数 \(\sim x^{m-n}\)，收敛当且仅当 \(m-n<-1\)，即 \(m<n-1\)。
\end{itemize}
因此，积分收敛当且仅当 \(m>-1\) 且 \(m<n-1\)。

(5) \(\int_{0}^{+\infty}\frac{\arctan x}{x^{n}}dx\)

解：需考虑 \(x\to0\) 和 \(x\to+\infty\) 的行为。
\begin{itemize}
\item 当 \(x\to0\) 时，\(\arctan x\sim x\)，被积函数 \(\sim\frac{x}{x^{n}}=x^{1-n}\)，收敛当且仅当 \(1-n>-1\)，即 \(n<2\)。
\item 当 \(x\to+\infty\) 时，\(\arctan x\to\frac{\pi}{2}\)，被积函数 \(\sim\frac{\pi/2}{x^{n}}\)，收敛当且仅当 \(n>1\)。
\end{itemize}
因此，积分收敛当且仅当 \(1<n<2\)。

(6) \(\int_{0}^{+\infty}\frac{\sin^{2}x}{x}dx\)

解：当 \(x\to0\) 时，\(\sin^{2}x\sim x^{2}\)，被积函数 \(\sim\frac{x^{2}}{x}=x\)，收敛。当 \(x\to+\infty\) 时，被积函数振荡，但绝对值行为类似于 \(\frac{1}{x}\)，而 \(\int\frac{1}{x}dx\) 发散，因此发散。故积分发散。

(7) \(\int_{0}^{\frac{\pi}{2}}\frac{dx}{\sin^{p}x\cos^{q}x}\)

解：奇点可能在 \(x=0\) 和 \(x=\frac{\pi}{2}\) 处。
\begin{itemize}
\item 在 \(x=0\) 处，\(\sin x\sim x\)，被积函数 \(\sim\frac{1}{x^{p}}\)，收敛当且仅当 \(p<1\)。
\item 在 \(x=\frac{\pi}{2}\) 处，\(\cos x\sim\frac{\pi}{2}-x\)，被积函数 \(\sim\frac{1}{\left(\frac{\pi}{2}-x\right)^{q}}\)，收敛当且仅当 \(q<1\)。
\end{itemize}
因此，积分收敛当且仅当 \(p<1\) 且 \(q<1\)。

(8) \(\int_{0}^{1}\frac{\ln x}{1-x^{2}}dx\)

解：奇点在 \(x=1\) 处。当 \(x\to1\) 时，\(1-x^{2}=(1-x)(1+x)\sim2(1-x)\)，\(\ln x\sim x-1\)，因此被积函数 \(\sim\frac{x-1}{2(1-x)}=-\frac{1}{2}\)，有限。在 \(x=0\) 处，\(\ln x\) 无界，但 \(\int\ln x\,dx\) 收敛。因此积分收敛。

6. 求下列极限

(1) \(\lim_{x\to0}x\int_{x}^{1}\frac{\cos t}{t^{2}}dt\)

解：使用 L'Hôpital 法则：
\[ \lim_{x\to0}x\int_{x}^{1}\frac{\cos t}{t^{2}}dt=\lim_{x\to0}\frac{\int_{x}^{1}\frac{\cos t}{t^{2}}dt}{1/x}. \]
求导：
\[ \frac{d}{dx}\int_{x}^{1}\frac{\cos t}{t^{2}}dt=-\frac{\cos x}{x^{2}},\qquad \frac{d}{dx}\frac{1}{x}=-\frac{1}{x^{2}}. \]
所以
\[ \lim_{x\to0}\frac{-\cos x/x^{2}}{-1/x^{2}}=\lim_{x\to0}\cos x=1. \]
故极限为 \(1\)。

(2) \(\lim_{x\to+\infty}\frac{\int_{0}^{x}\sqrt{1+t^{4}}\,dt}{x^{3}}\)

解：使用 L'Hôpital 法则：
\[ \lim_{x\to+\infty}\frac{\int_{0}^{x}\sqrt{1+t^{4}}dt}{x^{3}}=\lim_{x\to+\infty}\frac{\sqrt{1+x^{4}}}{3x^{2}}=\lim_{x\to+\infty}\frac{x^{2}\sqrt{1/x^{4}+1}}{3x^{2}}=\lim_{x\to+\infty}\frac{\sqrt{1+1/x^{4}}}{3}=\frac{1}{3}. \]

7. 设广义积分 \(\int_{1}^{+\infty}f^{2}(x)dx\) 收敛，证明广义积分 \(\int_{1}^{+\infty}\frac{f(x)}{x}dx\) 绝对收敛。

证明：要证明 \(\int_{1}^{+\infty}\frac{|f(x)|}{x}dx\) 收敛。由 Cauchy-Schwarz 不等式：
\[ \left(\int_{1}^{+\infty}\frac{|f(x)|}{x}dx\right)^{2}\le\int_{1}^{+\infty}f^{2}(x)dx\cdot\int_{1}^{+\infty}\frac{1}{x^{2}}dx. \]
已知 \(\int_{1}^{+\infty}f^{2}(x)dx\) 收敛，且 \(\int_{1}^{+\infty}\frac{1}{x^{2}}dx\) 收敛。因此，\(\left(\int_{1}^{+\infty}\frac{|f(x)|}{x}dx\right)^{2}\) 有限，故 \(\int_{1}^{+\infty}\frac{|f(x)|}{x}dx\) 收敛，即原积分绝对收敛。
